% ============================================================
%  chapters/ch08-semantic-renormalization.tex
% ============================================================

\chapter{Semantic Renormalization}

\chapterepigraph{%
  To understand is to forget the details\\
  while preserving the structure.\\
  Understanding is renormalization.
}{Flyxion, \textit{Semantic Infrastructure}}

\sidenote{Renormalization\\in physics:\\Wilson, 1971;\\ in logic:\\proof normalization}

The previous chapters treated evaluation as reduction: each pop step
removes one sphere, decreasing depth and complexity. This chapter examines
evaluation from the opposite direction — not as \emph{destruction} of
structure but as \emph{transformation} of structure. The relevant concept
is renormalization.

\section{The Collapse Operator Revisited}

In physics, renormalization is a procedure for handling theories with
infinitely many degrees of freedom by systematically integrating out
short-distance behavior and replacing it with effective parameters.
The key insight is that the result at long distances does not depend
on the microscopic details — only on certain aggregate properties
(coupling constants, anomalous dimensions) that survive the integration.

Spherepop evaluation has an identical structure.

\begin{definition}{Spherepop Renormalization}{sp-renorm}
  A \emph{Spherepop renormalization step} at depth $k$ collapses all
  spheres at depth $k$ simultaneously, replacing each with its evaluated
  value. The result is an expression of depth $k-1$ in which the
  depth-$k$ details have been integrated out.
\end{definition}

The depth-$k$ details (the specific sub-expressions that were collapsed)
are gone after renormalization. Only their effective values — the atoms
that replaced them — survive. The long-distance (low-depth) structure
of the expression is unaffected.

\section{Compression}

Renormalization produces \emph{compression}: a shorter, simpler
expression that computes the same thing as the original.

\begin{definition}{Semantic Compression}{sem-compress}
  An expression $E'$ is a \emph{semantic compression} of $E$ if:
  \begin{enumerate}
    \item $E' \equiv E$ (semantically equivalent: same terminal value),
    \item $D(E') < D(E)$ or $|E'| < |E|$ (strictly simpler).
  \end{enumerate}
\end{definition}

Every pop step is a semantic compression: the result has fewer spheres
and lower depth, but the same value. Full evaluation is maximal
compression: the terminal atomic sphere is the most compressed form
possible.

\begin{example}
  Starting from $E = \langle{+}\,:\, 1,\, \langle{\times}\,:\,3,\,
  \langle{\uparrow}\,:\,2,\,2\rangle\rangle\rangle$ (6 spheres, depth 2):
  \begin{align*}
    E &\to_\text{pop} \langle{+}\,:\, 1,\, \langle{\times}\,:\,3,\,4\rangle\rangle
      & \text{(5 spheres, depth 1)}\\
    &\to_\text{pop} \langle{+}\,:\, 1,\, 12\rangle
      & \text{(3 spheres, depth 1)}\\
    &\to_\text{pop} 13
      & \text{(1 sphere, depth 0)}
  \end{align*}
  Three compressions; final form is maximally compressed.
\end{example}

\section{Abstraction}

Renormalization in physics produces \emph{effective theories}: descriptions
valid at a given scale that replace the microscopic theory with a coarser
but computationally tractable model. Spherepop renormalization produces
the same: \emph{abstractions}.

\begin{definition}{Semantic Abstraction}{sem-abstract}
  A \emph{semantic abstraction} is a named partial expression:
  a sphere with a label $f$ that represents a complex sub-expression,
  without expanding the sub-expression. In bracket notation:
  $f$ (where $f$ is understood to expand to $\langle{+}\,:\,1,\,12\rangle$
  if needed).
\end{definition}

Abstraction is deferred renormalization: we give a name to the
result of a sub-computation without performing it. Functions, modules,
classes, and APIs are all forms of semantic abstraction — they provide
the \emph{interface} of a computation without revealing the
\emph{implementation}.

\begin{remark}
  The power of abstraction lies precisely in what it hides. A function
  $f(x) = x^2 + 2x + 1$ and a function $g(x) = (x+1)^2$ are
  semantically equivalent ($f \equiv g$) but structurally different.
  The abstracted form $h(x) = (x+1)^2$ is more compressed and reveals
  the structure (a perfect square) that the expanded form hides. Good
  abstractions choose the compressed representation that makes relevant
  structure visible while hiding irrelevant detail.
\end{remark}

\section{Quotient Structures}

Semantic renormalization produces quotient structures: the set of
all expressions, quotiented by semantic equivalence, organized by
their compression level.

\begin{definition}{Renormalization Tower}{renorm-tower}
  The \emph{renormalization tower} of an expression space $\mathcal{E}$
  is the sequence of quotient spaces:
  \[
    \mathcal{E} = \mathcal{E}_0 \twoheadrightarrow \mathcal{E}_1
    \twoheadrightarrow \cdots \twoheadrightarrow \mathcal{E}_n
  \]
  where $\mathcal{E}_k$ is the space of expressions of depth $\leq D-k$
  (expressions with the first $k$ levels of detail collapsed), and each
  arrow is a quotient projection.
\end{definition}

This tower is the Spherepop version of the renormalization group in
physics: a sequence of effective theories, each coarser than the last,
each valid at a different scale of detail.

\begin{figure}[h]
  \centering
  \begin{tikzpicture}[
    tower/.style={draw=RSVPdeep, thick, rounded corners=4pt,
      minimum width=5.5cm, minimum height=0.8cm,
      fill=RSVPlight, font=\small},
    arrowstyle/.style={-{Stealth[length=6pt]}, RSVPaccent, thick}
  ]
    \node[tower] (E0) at (0, 3.6)
      {$\mathcal{E}_0$: full expressions (depth $\leq D$)};
    \node[tower] (E1) at (0, 2.4)
      {$\mathcal{E}_1$: depth $\leq D{-}1$ (one level collapsed)};
    \node[tower] (E2) at (0, 1.2)
      {$\mathcal{E}_2$: depth $\leq D{-}2$ (two levels collapsed)};
    \node[tower, fill=RSVPdeep, text=white] (En) at (0, 0.0)
      {$\mathcal{E}_D$: atoms only (fully evaluated)};
    \draw[arrowstyle] (E0) -- (E1)
      node[midway, right=3pt, font=\footnotesize\itshape] {renormalize};
    \draw[arrowstyle] (E1) -- (E2)
      node[midway, right=3pt, font=\footnotesize\itshape] {renormalize};
    \draw[arrowstyle, dashed] (E2) -- (En)
      node[midway, right=3pt, font=\footnotesize\itshape] {$\cdots$};
  \end{tikzpicture}
  \caption{The renormalization tower of Spherepop expressions.
    Each level integrates out one more layer of syntactic detail,
    replacing it with effective values. The bottom level contains
    only atoms: fully evaluated, maximally compressed expressions.}
  \label{fig:renorm-tower}
\end{figure}

\section{Relationship to Admissibility}

Renormalization and admissibility interact through the constraint structure.
Each level of the renormalization tower has its own admissibility conditions:
what counts as a well-formed expression of depth $\leq k$.

\begin{proposition}{Admissibility is Renormalization-Stable}{adm-stable}
  If $E$ is admissible (satisfies all syntactic and semantic constraints),
  then any renormalization $R_k(E)$ is also admissible.
\end{proposition}

\begin{proof}
  Renormalization replaces each sub-expression with its value.
  Values are always well-typed atoms. Therefore the renormalized
  expression satisfies all typing constraints. Semantic constraints
  (e.g., the result is in the correct domain) are preserved because
  the value is the actual evaluated result.
\end{proof}

\section{Connection to CLIO and Semantic Fields}

The renormalization tower provides a natural interpretation for the
CLIO (Constraint-Leveraged Inference Optimization) framework
(Chapter~15, Appendix~H). CLIO uses constraint satisfaction to prune
the search space. Renormalization tells us why this works:

\begin{conceptaside}{Why Renormalization Enables Inference}
  A naive inference system must consider all possible continuations of
  a partial expression. This is exponential in the number of open positions.

  Renormalization provides a shortcut: instead of exploring all
  continuations in the full expression space $\mathcal{E}_0$, the system
  works in the renormalized space $\mathcal{E}_k$ where the fine details
  have been integrated out. The search space is exponentially smaller.

  This is why CLIO — and why human experts — can solve problems that
  brute-force search cannot. The expert has internalized the renormalization
  of their domain: they work at the level of abstracted patterns rather
  than raw details. The patterns are the fixed points of the renormalization
  flow: the structures that survive repeated compression.

  In terms of the RSVP fields: renormalized expressions correspond to
  states of low $\kappa$ (constraint violation) and high $S$ (accessibility).
  The renormalization flow is a gradient flow toward these states.
\end{conceptaside}

\begin{exercise}
  Apply two levels of renormalization to the expression
  $((2+3) \times (4+1)) + ((6-1) \times (1+1))$.
  What expression results? What detail has been integrated out?
\end{exercise}

\begin{exercise}
  Define ``semantic lossiness'' of a renormalization step as
  $L_k = d(E, R_k(E))$ (edit distance between the original and
  renormalized expression). Show that $L_k$ is monotonically
  non-decreasing in $k$ (more renormalization = more lossiness).
  When is $L_k = 0$?
\end{exercise}

\begin{exercise}[title={Algebraic}]
  The renormalization tower is a sequence of surjections
  $\mathcal{E}_0 \twoheadrightarrow \cdots \twoheadrightarrow \mathcal{E}_D$.
  Show that this sequence forms an inverse system. What is the
  inverse limit $\varprojlim \mathcal{E}_k$? (Hint: consider what
  the inverse limit means in terms of ``infinitely detailed'' expressions.)
\end{exercise}
